% ================= CHAPTER 4 =================
\chapter{HASIL DAN PEMBAHASAN}

\section{Hasil 1}
Sajikan data hasil penelitian yang relevan dengan bagian ini secara sistematis. Berikan deskripsi yang jelas terhadap data tersebut dan sertakan pembahasan yang mendalam untuk menjelaskan makna atau pola yang ditemukan dalam konteks penelitian.

\section{Hasil 2}
Penyajian data hasil penelitian dapat dibagi menjadi beberapa bagian sesuai dengan kebutuhan. Gunakan tabel, gambar, atau grafik untuk mempermudah pemahaman, kemudian diikuti dengan analisis dan interpretasi yang menghubungkannya dengan teori atau literatur terkait.

\section{Hasil 3}
Sesuaikan pembagian sub-bab hasil dengan keperluan dan seberapa banyak data yang diperoleh. Pastikan setiap temuan dibahas secara kritis untuk menunjukkan kontribusinya terhadap pencapaian tujuan penelitian yang telah ditetapkan.
